\begin{table}[htbp]
\centering
\small
\caption{Study cities, freight infrastructure and analysis sample.}
\label{tab:t1_cities}
\begin{tabular}{lllrrrl}
\toprule
City & State & Configuration & Rail goods sheds & IWT terminals & ICD/CFS & Analysis sample \\
\midrule
Patna & Bihar & Inland river port (NW-1), weak rail freight & 3 & 6 & 1 & 12,807 \\
Indore & Madhya Pradesh & Landlocked, peripheral ICD capacity & 5 & 0 & 5 & 48,884 \\
Varanasi & Uttar Pradesh & Inland river port, recent MMLP investment & 4 & 4 & 1 & 9,289 \\
Kochi & Kerala & Coastal major port, linear conurbation & 11 & 9 & 8 & 27,759 \\
Kolkata & West Bengal & Metropolitan port city, compact core & 5 & 5 & 6 & 24,852 \\
\bottomrule
\end{tabular}
\begin{tablenotes}\footnotesize
\item Terminal inventories include operational and under-development facilities; status is reported in Table S5. Indore has no inland waterway terminal within its district; see Section 3.3.
\end{tablenotes}
\end{table}
